% =============================================================================
% COPYRIGHT AND DISCLAIMER
% =============================================================================
% This file is provided as a convenience template for competition participants.
% By using it, you accept full responsibility for compliance with all applicable
% copyright, licensing, and usage rules of the Association for Computing
% Machinery (ACM) and of any LaTeX packages you load (including, without
% limitation, the "acmart" document class and related ACM materials).
%
% Consult ACM's official publication, template, and copyright policies before
% using this setup for any purpose other than this competition. Do not assume
% that this template may be reused, redistributed, or submitted elsewhere
% without independently verifying ACM's current rules.
%
% Use of the "acmart" class, the "sigconf" option, or ACM-style formatting in
% this template does NOT mean that this report, this competition, its organizers,
% or any resulting PDF is sponsored, endorsed, affiliated with, or published by
% ACM. ACM trademarks and branding remain the property of ACM.
%
% Template packages live in preamble.tex.
% Citations live in main.bib. Edit this file with your paper content.
% =============================================================================

\documentclass[sigconf]{acmart}
% =============================================================================
% Template setup — participants normally do not edit this file.
% Author content belongs in main.tex; citations belong in main.bib.
% \documentclass lives in main.tex so Overleaf does not treat this file as root.
% =============================================================================

% --- Competition / non-ACM-submission friendly rights metadata -------------
\setcopyright{none}
\settopmatter{printacmref=false}
\acmConference[Kaggle Competition]{Filament Segmentation Challenge}{2026}{\href{https://www.kaggle.com/competitions/filament-segmentation-2026}{kaggle.com/competitions/filament-segmentation-2026}}
\acmBooktitle{Competition Report}
\acmDOI{}
\acmISBN{}

% --- Author guidance text --------------------------------------------------
% \guide{...} is gray italic while \guideStyletrue; omitted from the PDF if false.
% Flip the switch in main.tex after % =============================================================================
% Template setup — participants normally do not edit this file.
% Author content belongs in main.tex; citations belong in main.bib.
% \documentclass lives in main.tex so Overleaf does not treat this file as root.
% =============================================================================

% --- Competition / non-ACM-submission friendly rights metadata -------------
\setcopyright{none}
\settopmatter{printacmref=false}
\acmConference[Kaggle Competition]{Filament Segmentation Challenge}{2026}{\href{https://www.kaggle.com/competitions/filament-segmentation-2026}{kaggle.com/competitions/filament-segmentation-2026}}
\acmBooktitle{Competition Report}
\acmDOI{}
\acmISBN{}

% --- Author guidance text --------------------------------------------------
% \guide{...} is gray italic while \guideStyletrue; omitted from the PDF if false.
% Flip the switch in main.tex after % =============================================================================
% Template setup — participants normally do not edit this file.
% Author content belongs in main.tex; citations belong in main.bib.
% \documentclass lives in main.tex so Overleaf does not treat this file as root.
% =============================================================================

% --- Competition / non-ACM-submission friendly rights metadata -------------
\setcopyright{none}
\settopmatter{printacmref=false}
\acmConference[Kaggle Competition]{Filament Segmentation Challenge}{2026}{\href{https://www.kaggle.com/competitions/filament-segmentation-2026}{kaggle.com/competitions/filament-segmentation-2026}}
\acmBooktitle{Competition Report}
\acmDOI{}
\acmISBN{}

% --- Author guidance text --------------------------------------------------
% \guide{...} is gray italic while \guideStyletrue; omitted from the PDF if false.
% Flip the switch in main.tex after \input{preamble}.
%
% Do not use \guide{...} for \email{...}.
\usepackage{xcolor}
\newif\ifguideStyle
\guideStyletrue
\newcommand{\guide}[1]{%
  \ifguideStyle
    {\color{gray!65}\itshape #1}%
  \fi
}

% --- Math ------------------------------------------------------------------
% Do NOT load amssymb (or amsfonts): acmart already loads newtxmath, which
% defines the same symbols (e.g. \Bbbk) and causes a conflict on Overleaf.
% amsmath is already provided by acmart; mathtools extends it.
\usepackage{mathtools}

% --- Figures ---------------------------------------------------------------
% graphicx is already loaded by acmart
\usepackage{subcaption}

% --- Tables ----------------------------------------------------------------
% booktabs is already loaded by acmart
\usepackage{tabularx}
\usepackage{multirow}

% --- Bibliography style is set in main.tex (ACM-Reference-Format) ----------
.
%
% Do not use \guide{...} for \email{...}.
\usepackage{xcolor}
\newif\ifguideStyle
\guideStyletrue
\newcommand{\guide}[1]{%
  \ifguideStyle
    {\color{gray!65}\itshape #1}%
  \fi
}

% --- Math ------------------------------------------------------------------
% Do NOT load amssymb (or amsfonts): acmart already loads newtxmath, which
% defines the same symbols (e.g. \Bbbk) and causes a conflict on Overleaf.
% amsmath is already provided by acmart; mathtools extends it.
\usepackage{mathtools}

% --- Figures ---------------------------------------------------------------
% graphicx is already loaded by acmart
\usepackage{subcaption}

% --- Tables ----------------------------------------------------------------
% booktabs is already loaded by acmart
\usepackage{tabularx}
\usepackage{multirow}

% --- Bibliography style is set in main.tex (ACM-Reference-Format) ----------
.
%
% Do not use \guide{...} for \email{...}.
\usepackage{xcolor}
\newif\ifguideStyle
\guideStyletrue
\newcommand{\guide}[1]{%
  \ifguideStyle
    {\color{gray!65}\itshape #1}%
  \fi
}

% --- Math ------------------------------------------------------------------
% Do NOT load amssymb (or amsfonts): acmart already loads newtxmath, which
% defines the same symbols (e.g. \Bbbk) and causes a conflict on Overleaf.
% amsmath is already provided by acmart; mathtools extends it.
\usepackage{mathtools}

% --- Figures ---------------------------------------------------------------
% graphicx is already loaded by acmart
\usepackage{subcaption}

% --- Tables ----------------------------------------------------------------
% booktabs is already loaded by acmart
\usepackage{tabularx}
\usepackage{multirow}

% --- Bibliography style is set in main.tex (ACM-Reference-Format) ----------


% Guide styling: \guideStyletrue = show gray italic hints; \guideStylefalse = omit them
% \guideStyletrue
\guideStylefalse


\begin{document}

\title{Native-Resolution Mask R-CNN Ensemble with Flip Test-Time Augmentation}

\subtitle{A Solution to the Solar Filament Segmentation Challenge 2026} % DO NOT CHANGE


\author{Shamsul Hakimee}
\affiliation{%
  \department{Kulliyyah of Engineering}
  \institution{International Islamic University Malaysia}
  \city{Kuala Lumpur}
  \country{Malaysia}
}
\email{shamsul.hakimee@student.iium.edu.my}

\author{Amirul Iman}
\affiliation{%
  \department{Kulliyyah of Engineering}
  \institution{International Islamic University Malaysia}
  \city{Kuala Lumpur}
  \country{Malaysia}
}
\email{amiruliman.khir@student.iium.edu.my}

\author{Syazani Akmal}
\affiliation{%
  \department{Kulliyyah of Engineering}
  \institution{International Islamic University Malaysia}
  \city{Kuala Lumpur}
  \country{Malaysia}
}
\email{syazani.akmal@student.iium.edu.my}


\begin{abstract}
    This report describes our solution to the Solar Filament Segmentation Challenge~2026, a Kaggle competition on automatic segmentation of solar filaments in GONG H-$\alpha$ observations.
Our pipeline uses two Mask R-CNN models (ResNeXt-101 and ResNet-50 FPN
backbones) that work on the full $2048\times2048$ images without tiling,
with anchors shaped for thin, elongated filaments. Each model predicts
four flipped views of every image; the views are merged per model and the
two models are then merged by mask intersection, followed by simple
clean-up rules that target the fragmentation and false-alarm penalties
of panoptic quality (PQ). Every annotator's labels are used as separate
training samples. On a held-out validation split built from 27-day time
blocks, the final pipeline reaches PQ~0.448 (95\% CI 0.425--0.470) and a
mean Dice of 0.652, against an inter-annotator agreement of PQ~0.336;
it scores 0.39 on the public leaderboard. Code, weights and a one-click
notebook are public.
% source: results/phase8/Mirul/rubric_exp101.json (PQ, bootstrap CI, Dice_samplewise); README.md (LB 0.39); docs/HANDOVER.md §5 (0.336)
\end{abstract}


\maketitle

\ifguideStyle
\noindent{\color{gray!65}\itshape
Guide mode is on.}
\par\medskip
\fi

\guide{Please limit your report to four (4) pages, excluding the Acknowledgment and References sections. Throughout this document, please follow these considerations: 
\begin{itemize}
    \item Do not alter any text that is currently in black color (except for the emails under authors' name).
    \item Replace all gray text with your own.
    \item When you are done with the report, please replace \texttt{\textbackslash guideStyletrue} with \texttt{\textbackslash guideStylefalse}, to remove any gray text. Alternatively, you can remove any text wrapped in \texttt{\textbackslash guide\{\}},
    \item Be as succinct as possible, and only provide details regarding your own work. For example, there is no need to discuss what solar filaments are, or why filament segmentation matters.
    \item If you must, you can add extra sections.
\end{itemize}}


% ==================================== %
%              INTRODUCTION            %
% ==================================== %
\section{Introduction}
    This report goes over the details of our submitted solution to the Solar Filament Segmentation Challenge~2026 \cite{filament-segmentation-2026}. The challenge, launched on July~10,~2026, is organized by the Earth-Space AI Research (ESAIR) Lab\footnote{\href{https://www.esairlab.com/home}{www.esairlab.com/}}. This challenge utilizes the MAGFiLO dataset for evaluation of the segmentation algorithms \cite{ahmadzadeh2024dataset}. The task is instance segmentation: every filament must be a separate
mask. Submissions are scored with panoptic quality,
\begin{equation}
  \mathrm{PQ}=\underbrace{\frac{\sum_{(p,g)\in TP}\mathrm{IoU}(p,g)}{|TP|}}_{\mathrm{SQ}}
  \times\underbrace{\frac{|TP|}{|TP|+\frac12|FP|+\frac12|FN|}}_{\mathrm{RQ}},
\end{equation}
where a prediction and a ground-truth mask form a true positive (TP)
only if their IoU is above 0.5. Two properties of the data make this hard.
First, filaments are long and thin, so a prediction that breaks one
filament into pieces gains one TP but adds false positives (FP). Second,
the labels are noisy: when two annotators outline the same image, their
masks agree only at PQ~$\approx$~0.336.
% source: SelfEvaluationResult/FINDINGS.md (official metric, IoU > 0.5); docs/HANDOVER.md §5 (0.336)

Our solution rests on three ideas, each kept only if it beat the
validation noise of about 0.010~PQ (Sec.~\ref{sec:evaluation}):
(i) train and predict at the native $2048\times2048$ resolution, so thin
filaments are not lost by down-scaling; (ii) combine two different
Mask R-CNN models and keep only the pixels both agree on; (iii) predict
each image in four flipped views and average them, which lowers the
scores of chance detections. Sec.~\ref{sec:methodology} describes the
data, the models and the pipeline; Sec.~\ref{sec:evaluation} reports
the results, what did not work, and the remaining errors.
% source: README.md "How we got here"; experiments/ablation_table.md (Phase 6 M1 noise threshold 0.010)



% ==================================== %
%              METHODOLOGY             %
% ==================================== %
\section{Methodology}\label{sec:methodology}
% ============================================================
% 2a  Data and validation split   (Syaz, S8.1)
% ============================================================
\subsection{Data and validation split}\label{sec:data}
We use only the competition data: full-disk GONG H$\alpha$ images from
the MAGFiLO~1.0 dataset~\cite{ahmadzadeh2024dataset}, 707 annotated training files and
180 unannotated test files, all $2048\times2048$ JPEGs.
% source: CLAUDE.md §2 (dataset layout); docs/HANDOVER.md §5
The training annotations hold 8{,}199 filament outlines, with areas from
9 to 37{,}739~px$^2$ (mean $\approx$2{,}120~px$^2$).
% source: CLAUDE.md §2 (instance count and area range)

\textbf{Several annotators per image.} The annotation file has 1{,}154
image entries for the 707 files: 296 files were outlined independently by
two or three annotators.
% source: docs/HANDOVER.md §5; docs/COMPLETED_TASKS.md Phase 1
The annotators often disagree. Scored against each other with the
competition metric, two annotators of the same image reach only
PQ~$\approx$~0.336 (SQ~0.63, RQ~0.53), a rough human ceiling for the task.
% source: docs/COMPLETED_TASKS.md Phase 1; docs/HANDOVER.md §5
We never merge the annotation sets. For training, every annotator set is
its own sample (921 samples from 566 training files); for evaluation, all
sets of a file are used, as in the official scorer.
% source: README.md "Training"; docs/HANDOVER.md §5; SelfEvaluationResult/FINDINGS.md

\textbf{Validation split.} We hold out 141 of the 707 files (20\%) as a
validation set. It was frozen once (seed 42) and used for every model
choice.
% source: splits/val_ids.json (n_train 566, n_val 141, seed 42)
Images taken a few days apart show the same filaments, so a random split
would leak near-copies into validation. We therefore group the files into
27-day blocks (about one solar rotation) and keep each block entirely in
training or in validation (stratified group 5-fold, strata = number of
filaments per file; mean 7.46 in training vs.\ 7.31 in validation).
% source: splits/val_ids.json ("method", train_mean_filaments, val_mean_filaments)
The test set has no ground truth and was used only for leaderboard
submissions. All validation numbers in this report are the official PQ on
these 141 files; the public leaderboard (LB) uses about half of the test
images.
% source: docs/HANDOVER.md §2b "Rules for the report"

% ============================================================
% 2b  Models and training   (Syaz, S8.2)
% ============================================================
\subsection{Models and training}\label{sec:models}
We use two Mask R-CNN instance-segmentation models from Detectron2, both
pretrained on COCO: one with a ResNeXt-101 (32$\times$8d) FPN backbone and
one with a ResNet-50 FPN backbone. Mask R-CNN first proposes boxes around
possible filaments, then draws a mask inside each box and gives it a
confidence score. We use a single class (``filament'').
% source: src/config.py (MODEL_CONFIGS, NUM_CLASSES); src/final_pipeline.py header
Three changes adapt the models to filaments:
(i) small anchor boxes of 8--128~px with aspect ratios 0.2--5, because
filaments are long and thin;
% source: src/config.py ModelConfig (ANCHOR_SIZES, ASPECT_RATIOS)
(ii) for the X101 model, a finer mask head that draws each mask on a
$56\times56$ grid instead of the default $28\times28$, which keeps thin
filament edges;
% source: src/config.py (MASK_POOLER_RES 14 -> 28); src/final_pipeline.py MODELS
(iii) both models train and predict on the full $2048\times2048$ image,
without tiling or resizing, so thin filaments keep their fine detail.
No pre-processing is applied: solar-disk masking and CLAHE contrast
enhancement gave no gain (Sec.~\ref{sec:evaluation}).
% source: README.md "The final pipeline" (Input row) and "How we got here"; experiments/ablation_table.md

Both models use the same training recipe: every annotator set of the 566
training files is its own sample (921 samples); 6{,}000 iterations with
batch size 2; SGD with learning rate 0.0025, a 200-iteration linear
warm-up and a 10$\times$ drop at 60\% and 85\% of training; random
horizontal and vertical flips and brightness and contrast changes of
$\pm$20\%; seed 42. Training ran on one NVIDIA A100 GPU (Google Colab).
% source: README.md "Training"; src/config.py TrainConfig; src/train.py (WARMUP_ITERS, STEPS, GAMMA); src/augment.py (RandomFlip, RandomBrightness/Contrast 0.8-1.2)

% ============================================================
% 2c  Inference, fusion and post-processing   (Syaz, S8.2)
% ============================================================
\begin{figure*}[t]
  \centering
  \includegraphics[width=\textwidth]{fig_pipeline.pdf}
  \caption{The final pipeline. Each test image is predicted by two Mask R-CNN
  models on four flipped views; the views of each model are merged first
  (Fusion~1), then the two models (Fusion~2), and the result is cleaned up
  before it is written to the submission file.}
  \Description{Block diagram: test image, four flipped views, two Mask R-CNN
  models, two view-fusion steps, model fusion, clean-up, submission CSV.}
  \label{fig:pipeline}
\end{figure*}
% source: docs/report/make_fig_pipeline.py (settings from src/final_pipeline.py)

\subsection{Inference, fusion and post-processing}\label{sec:inference}
The full pipeline is shown in Fig.~\ref{fig:pipeline}.
For each test image, each model predicts four views: the original image
and its horizontal, vertical and double flip. The predicted masks are
flipped back, and only masks with a score of at least 0.2 are kept at
this stage.
% source: src/final_pipeline.py (VIEWS; DUMP_SCORE_THRESHOLD 0.2 in src/config.py)
The predictions are then merged in two steps:
\begin{enumerate}
  \item \textbf{Views (per model).} Masks from different views with
  IoU~$\geq$~0.5 are the same filament. The merged mask is their union and
  its score is the mean over the four views; a view that missed the
  filament counts as 0. A filament seen in only one view therefore gets a
  low score.
  \item \textbf{Models.} The two models are merged the same way, with
  IoU~$\geq$~0.3; the merged mask is the intersection of the two masks and
  its score is the mean of the two models.
\end{enumerate}
% source: src/final_pipeline.py header and VIEW_IOU/VIEW_MODE/MODEL_IOU/MODEL_MODE
Finally, four post-processing steps target the metric's failure modes:
(i) morphological closing with a 5-px elliptical kernel reconnects small
breaks inside a filament, which would otherwise count as extra false
positives; (ii) masks with a score below 0.6 are removed; (iii) masks
smaller than 50~px$^2$ are removed as noise; (iv) overlapping masks are
resolved: the higher-scored mask keeps the shared pixels, and a mask that
loses more than 20\% of its area is dropped. The remaining masks are
written as COCO run-length encodings in the submission file.
% source: src/final_pipeline.py (SCORE_THR 0.6, CLOSE_KERNEL 5, MIN_AREA 50, DROP_FRAC 0.2); README.md "The final pipeline"
Fig.~\ref{fig:examples} shows the output on three validation images.

\begin{figure*}[t]
  \centering
  \includegraphics[width=\textwidth]{fig_examples.pdf}
  \caption{Example segmentations on validation images, chosen by their own
  PQ: (a) a good image (90th percentile), (b) a typical one (median) and
  (c) a hard one (10th percentile). For each image, the left crop shows one
  annotator's filaments and the right crop shows our final masks judged
  against that annotator. In (c), a mask that covers a real filament but
  overlaps it with IoU below 0.5 counts as both a false alarm and a miss.}
  \Description{Six image crops of the Sun in H-alpha in three pairs. In each
  pair, annotator outlines on the left and our coloured masks on the right:
  green hits, red false alarms, dashed cyan missed filaments.}
  \label{fig:examples}
\end{figure*}
% source: docs/report/make_fig_examples.py (val images 20111112193254Bh,
% 20130707134154Bh, 20110923165814Mh; image PQ 0.56 / 0.45 / 0.32)
% ============================================================
% 2d  Implementation and reproducibility   (Syaz, S8.5)
% ============================================================
\subsection{Implementation and reproducibility}\label{sec:code}
All code, the frozen validation split and the trained weights of both
models are in our public repository\footnote{\url{https://github.com/amiruliman592/Kaggle-Competition-}};
nothing else has to be downloaded. The weights are stored in 94-MB parts
with SHA-256 checksums and joined by one command. The notebook
\texttt{full\_\allowbreak pipeline.ipynb} (folder \texttt{notebooks/})
runs the whole pipeline on a Colab
GPU (``Run all''): it predicts the 180 test images, writes and validates
the submission file and scores the validation split. A run from a fresh
clone gave the same validation result (PQ 0.4480, TP/FP/FN 1083/495/595)
and found 1{,}255 of the 1{,}256 submitted test masks again. On an NVIDIA L4
the whole pipeline takes 6.0~s per image (median; 8.2~s at the 95th
percentile), 18 minutes for the 180 test images, with 2.9~GB peak GPU
memory. The repository also holds 84 automated tests, a README guide for
reviewers and a plain-language notebook that explains every step
(\texttt{notebooks/pipeline\_story.ipynb}).
% source: README.md "Reproduce the final submission"; results/phase7/final_run/report_test.json (median 6.02 s, p95 8.21 s, total 1083 s, 2.94 GB);
% results/phase7/final_run/report_val.json; docs/COMPLETED_TASKS.md Phase 7 (clean clone, 1,255 of 1,256); weights/manifest.json; tests/ (84 test functions)




    \guide{Explain your methodology, ensuring that the following items are included:
    \begin{itemize}
        \item (for novel algorithms) a detailed description of the algorithm introduced for segmentation,
        \item (for existing algorithms) a quick review of the algorithm used for segmentation, and appropriate citation(s),
        \item a detailed description of the end-to-end data pipeline, including pre-processing and post-processing,
        \item (for ML models) a detailed description of any sampling strategy including train-val-test splits,
        \item a graphic illustrating the entire data pipeline,
        \item one or more graphics showing a few examples of the segmentations.
    \end{itemize}}

% ==================================== %
%              EVALUATION              %
% ==================================== %
\section{Evaluation}\label{sec:evaluation}
% ============================================================================
% sec_evaluation.tex  --  Section 3 Evaluation (Mirul / Amirul Iman), Phase 8 M8.2-M8.7
% \input under "Evaluation" in the Overleaf main file. Needs fig_dice_iou.pdf uploaded.
% Every number has its source file in a comment. Parts marked "% OPTIONAL" can be cut
% first if the report is over 4 pages.
% ============================================================================

% ---------------------------------------------------------------- M8.2
\subsection{Evaluation protocol}
All choices were made on a frozen validation split of 141 of the 707 training files (27-day time blocks, so no
validation day is also a training day). We score with the official Kaggle PQ: every annotator set of a file is its
own ground-truth entry (233 entries), a prediction matches a ground-truth filament at IoU $>0.5$, and TP, FP and FN
are pooled over all entries.
% results/phase8/Mirul/rubric_exp101.json (num_entries 233)
Our implementation gives the same value as the organisers' notebook on the same files (0.3296 both).
% SelfEvaluationResult/FINDINGS.md
Dice is \texttt{torchmetrics} \texttt{DiceScore} on the binary foreground mask (union of all filaments), averaged
over entries; the organisers' exact settings are not published.

The labels are noisy: two annotators of the same image reach only PQ 0.336 against each other, and an exact copy of
one annotator scores 0.739 on our validation entries, so the reachable PQ is far below 1.
% output/eda/eda_summary.json (0.336); SelfEvaluationResult/FINDINGS.md (0.7394)
To separate real gains from chance we used a paired bootstrap over the 141 files (2,000 draws): a gain smaller than
about 0.010 PQ is not distinguishable from noise, and retraining with another random seed changed PQ by only 0.001.
% results/phase6/Mirul/bootstrap.json; experiments/ablation_table.md Phase 6 M1a, S1
The public leaderboard (about half of the 180 test images, two decimals) has a standard deviation of about 0.014,
so we used it only as a check. Validation PQ was 0.05--0.06 above the leaderboard on every upload, but the ranking
of the steps in Table~\ref{tab:journey} was the same on both.
% experiments/ablation_table.md Phase 6 M1a (LB SD 0.013-0.015); experiments/log.csv (lb_score column)

% ---------------------------------------------------------------- M8.3 + M8.1
\subsection{Results}
Table~\ref{tab:journey} shows how each step changed the score. The largest gains came from native 2048-pixel input
($+0.025$), from combining two models ($+0.010$) and from the four flipped views with two-stage fusion ($+0.011$).
The last two are above the noise level: $+0.039$ [0.027, 0.050] for the ensemble against the 1024-pixel model and
$+0.012$ [0.004, 0.020] for the final pipeline against the plain ensemble (95\% bootstrap intervals).
% results/phase6/Mirul/bootstrap.json (exp010 vs exp006); results/phase6/Mirul/bootstrap_m2.json (exp101 vs exp010)

\begin{table}[t]
\centering
\caption{From baseline to final pipeline. Validation: official PQ and mean Dice on the frozen 141-file split.
LB: public leaderboard, about 50\% of the test set.}
\label{tab:journey}
\small
\begin{tabular}{@{}lccc@{}}
\hline
Step & Val PQ & Dice & LB \\
\hline
Mask R-CNN R50, 1024 px (baseline)     & 0.391 & 0.609 & 0.34 \\ % log.csv exp003
+ all annotator sets as training data  & 0.398 & 0.643 & 0.35 \\ % log.csv exp006
+ native 2048 px, R50                  & 0.423 & 0.648 & 0.36 \\ % log.csv exp007 (thr 0.6)
\quad X101 instead of R50              & 0.426 & 0.642 & 0.37 \\ % log.csv exp008
+ ensemble X101 + R50                  & 0.437 & 0.651 & 0.38 \\ % log.csv exp010
+ 56$\times$56 mask head, 4 flip views, & \textbf{0.448} & \textbf{0.652} & \textbf{0.39} \\ % log.csv exp101
\quad two-stage fusion (\textbf{final}) & & & \\
\hline
\end{tabular}
\end{table}

\textbf{Final pipeline.} On validation it reaches PQ $0.448$ (95\% interval $0.425$--$0.470$), with segmentation
quality SQ $0.673$ and recognition quality RQ $0.665$ (1,083 TP, 495 FP, 595 FN), and mean Dice $0.652$
($0.628$--$0.672$). The intervals come from 2,000 bootstrap draws of the 141 files.
% results/phase8/Mirul/rubric_exp101.json (PQ, SQ, RQ, TP/FP/FN, Dice_samplewise, bootstrap_95ci_files)
Fig.~\ref{fig:diceiou} shows the Dice and IoU of every pair of a predicted and a ground-truth filament that overlap.
The matched pairs are tightly grouped (median Dice 0.81, median IoU 0.67); the long tail below the 0.5 cut-off are
filaments we found but outlined too differently from the annotator.
% rubric_exp101.json dist_Dice_matched_TP.median 0.806, dist_IoU_matched_TP.median 0.675
\textbf{One-to-many and many-to-one.} Of the ground-truth filaments, 20 were split into several predicted pieces
(19 into two, one into three); 58 of our 933 masks cover two or more ground-truth filaments (55 cover two, 3 cover
three). A piece counts when at least half of its area lies inside the other mask. Merging is therefore about three
times more common than fragmentation.
% rubric_exp101.json one_to_many_*, many_to_one_*, n_pred_masks 933

\begin{figure}[t]
\centering
\includegraphics[width=\columnwidth]{fig_dice_iou.pdf}
\caption{Dice (left) and IoU (right) of every overlapping pair of predicted and ground-truth filaments on
validation (all annotator sets). Dark: pairs counted as true positives (IoU $>0.5$, dashed line).}
\label{fig:diceiou}
% docs/report/m81_rubric.py -> docs/report/figures/fig_dice_iou.pdf
\end{figure}

% ---------------------------------------------------------------- M8.4
\subsection{What did not work}
Table~\ref{tab:negative} lists ideas that did not pass the 0.010 noise bar. Training-recipe changes were tested on
the cheaper 1024-pixel R50 model (reference PQ 0.398), all others on the 2048-pixel models.

\begin{table}[t]
\centering
\caption{Rejected ideas (validation official PQ; $\Delta$ against the reference of that test).}
\label{tab:negative}
\small
\begin{tabular}{@{}lcl@{}}
\hline
Idea & $\Delta$PQ & Reason \\
\hline
CLAHE + solar-disk mask      & $+0.000$ & no gain \\                    % ablation_table Retraining A
SAM 2 mask refinement        & $-0.019$ & tighter than annotators \\    % ablation_table Phase 5 step 1
Mask2Former (Swin-S)         & $-0.024$ & twice the merges \\           % ablation_table Phase 5 step 2
3$\times$ longer schedule    & $+0.001$ & noise, lower Dice \\          % ablation_table S2 (exp202)
90$^\circ$ rotations + zoom  & $-0.001$ & no gain \\                    % ablation_table S3 (exp203)
Copy-paste of filaments      & $-0.004$ & more false alarms \\          % ablation_table S4 (exp204)
Self-training, 2,098 images  & $-0.002$ & model too cautious \\         % ablation_table S5 (exp205)
Third ensemble member        & $-0.003$ & same errors as X101 \\        % ablation_table Phase 7 prep A
Learned re-scoring           & $-0.006$ & too few labels \\             % ablation_table Phase 7 prep D (exp016)
Retrain on all 707 files     & LB $=$   & 95\% same masks \\            % ablation_table Phase 7 prep B (exp015)
\hline
\end{tabular}
\end{table}

% OPTIONAL (cut first): one-sentence explanations of Table 2
SAM 2 draws the dark core of a filament, about 15\% smaller than the human outline, which confirms that the labels,
not mask sharpness, limit SQ. Self-training used only confident pseudo-labels (score $\geq 0.85$), so faint
filaments in the extra images were taught as background.
% ablation_table Phase 5 step 1, Phase 6 S5

% ---------------------------------------------------------------- M8.5
\subsection{Error analysis}
We put every validation error of the final pipeline into one cause and computed the PQ if only that cause were
fixed. The largest is confidence: 282 ground-truth filaments are found by a raw prediction (IoU $>0.5$) but scored
below the 0.6 cut-off ($+0.072$ PQ if kept; 233 of them medium-sized). Next are 305 false alarms that no annotator
drew ($+0.046$), 162 filaments with no prediction near them ($+0.043$) and 115 near misses with IoU 0.3--0.5
($+0.031$). A further 190 ``false alarms'' were drawn by another annotator of the same image, i.e.\ label noise
rather than model error.
% results/phase6/Mirul/error_budget_exp101.json; experiments/ablation_table.md Phase 6 M4
By size, medium filaments score PQ 0.430 against 0.493 for large ones; the 12 small filaments ($<200$\,px$^2$)
are all missed, which costs at most 0.005 PQ. By station, PQ ranges from 0.418 (El Teide, most false alarms) to
0.468 (Big Bear); with only 20--28 files per station these gaps are within noise.
% error_budget_exp101.json by size / station; docs/HANDOVER.md section 6 (small <= 0.005)

% ---------------------------------------------------------------- M8.6
\subsection{Runtime}
Inference runs eight forward passes per image (two models $\times$ four views) plus fusion and post-processing:
median 6.0\,s and 95th percentile 8.2\,s per 2048$\times$2048 image on one NVIDIA L4 GPU, 18 minutes for all 180
test images, with 2.9\,GB peak GPU memory. Training at 2048 pixels needs about 20\,GB and was done on an A100.
% results/phase7/final_run/report_test.json; experiments/log.csv exp008 (peak 20.2 GB)

% ---------------------------------------------------------------- M8.7
\subsection{Strengths, weaknesses and improvements}
\textbf{Strengths.} Every step was chosen on a fixed split against a measured noise level, so the gains carried
over to the leaderboard (0.34 to 0.39). The method is simple (two standard Mask R-CNN models, no pre-processing) and
a clean clone of the repository reproduces the validation result exactly.
% results/phase7/final_run/report_val.json (PQ 0.4480 again)
\textbf{Weaknesses.} The main loss is confidence near the cut-off: many correct medium-sized filaments are scored
just below 0.6, and one threshold cannot keep them without adding false alarms. Over-merging of nearby filaments is
three times more common than fragmentation, and small filaments are not found.
\textbf{Improvements.} Better calibrated scores for medium filaments, for example from more labelled data or
self-training with a lower pseudo-label threshold, which we did not have time to test; and a rule or model to
split masks that join two filaments.
% error_budget_exp101.json (fn_low_score); ablation_table Phase 6 S5 (lower threshold not tried)

    \guide{Explain your evaluation. This includes:
    \begin{itemize}
        \item the evaluation strategy,
        \item plots and/or tables showing the performance of the segmentation algorithm (please include error bars and uncertainty measures whenever possible),
        \item strengths and weaknesses of the proposed algorithm, and possible strategies for improvement.
    \end{itemize}}



% ==================================== %
%            ACKNOWLEDGMENT            %
% ==================================== %
\section{Acknowledgment}\label{sec:acknowledgment}
    % Do not change
    The organization of the Filament Segmentation Challenge 2026 is partially supported by the U.S. National Science Foundation (NSF) under Grant No. 2209912 and 2433781, directorate for Computer and Information Science and Engineering (CSE), and Office of Advanced Cyberinfrastructure (OAC), and Grant No. 2511630, AST Division Of Astronomical Sciences and MPS Directorate for Mathematical and Physical Sciences.

    The Filament Segmentation Challenge 2026 is sponsored by the U.S. National Science Foundation (NSF) National Solar Observatory (NSO). 

    This work utilizes GONG data obtained by the NSO Integrated Synoptic Program, managed by the National Solar Observatory, which is operated by the Association of Universities for Research in Astronomy (AURA), Inc. under a cooperative agreement with the National Science Foundation and with contribution from the National Oceanic and Atmospheric Administration. The GONG network of instruments is hosted by the Big Bear Solar Observatory, High Altitude Observatory, Learmonth Solar Observatory, Udaipur Solar Observatory, Instituto de Astrofísica de Canarias, and Cerro Tololo Interamerican Observatory.

\bibliographystyle{ACM-Reference-Format}
\bibliography{main}

\end{document}
